\chapter*{Main theorem}
\addcontentsline{toc}{chapter}{Main theorem}

The construction starts from Wilson lattice gauge theory on periodic
four-dimensional lattices. Its continuum observables are local
gauge-invariant fields. Their Euclidean correlations determine the
relativistic theory described below.

\noindent\textbf{Main theorem (continuum Yang--Mills fields).}
Let $G$ be a compact matrix group with simple nonabelian real Lie algebra
$\gLie$. Fix a faithful unitary representation and its normalized trace
inner product, a positive penalty parameter for the block transformation,
and a sufficiently large inverse coupling at the coarsest scale.
There is a choice of lattice spacings tending to zero, admissible physical
periods tending to infinity, bare Wilson couplings, and normalizations of
local observables with the following properties.

\begin{enumerate}[label=\textup{(\roman*)},leftmargin=*]
\item \textit{Continuum limit.} At each fixed admissible period the
correlations converge along the full sequence of lattice spacings. From
any increasing sequence of these periods one can extract a single
subsequence along which the correlations of all orders converge to smooth
functions at pairwise distinct points in $\RR^4$. The same subsequence
works for all derivatives. The resulting local fields form nested
finite-dimensional spaces, obtained by imposing successive bounds on the
scaling dimension of the polynomial labels (curvature counts as two, each
derivative as one, and dimensions add in products). They contain the identity and
admit commuting spacetime derivatives. At each bound the field space and
its gauge-invariant free Maxwell counterpart have the same dimension and
carry isomorphic representations of $O(4)$.

\item \textit{Relativistic theory.} The limiting correlations are averages,
with respect to a probability measure of full support on a compact metrizable
space of phases, of Euclidean covariant, reflection-positive correlations.
Every phase has the same field spaces and reconstructs a local,
Poincar\'e-covariant theory on four-dimensional Minkowski space. Its
Wightman distributions are tempered, its energy--momentum spectrum lies
in the closed forward light cone, and its invariant vacuum is unique up
to scalar multiples. The weak commutant of the fields consists of scalars.

\item \textit{Stress tensor and local products.} The phases share a real,
symmetric, conserved stress tensor whose local charges generate translations
and rotations through the Ward identities. Products of nearby fields admit
an operator product expansion to each fixed finite order, with coefficient
and remainder bounds uniform in the phase. A fixed Taylor subtraction
also extends the separated Euclidean correlations across coincident points
as distributions of locally finite order.
\end{enumerate}

Vacuum uniqueness here is a statement within each phase. Uniqueness of
the infinite-volume limit and a positive mass gap remain open in this construction.

\clearpage
\section*{Where the theorem is proved}
\addcontentsline{toc}{section}{Where the theorem is proved}

Theorem~\ref{ch19:conditional} (p.~\pageref{ch19:conditional}) assembles
the conclusions for one and the same family of Wilson measures. The main
steps of its proof can be read separately as follows.

\paragraph{The two limits in (i).}
Theorems~\ref{ch10:finite-history} and~\ref{ch10:convergence} tune the
bare coupling and prove convergence of the renormalization group trajectory
at each fixed scale. With the bounds on observable insertions from
Chapter~\ref{ch11:insertions}, Theorem~\ref{ch12:uv-theorem} proves the
full cutoff limit at fixed volume. Section~\ref{ch13:extraction} then
extracts one subsequence of volumes for all separated correlations and
their derivatives. Proposition~\ref{ch13:rank} identifies the dimensions
of the field spaces. Section~\ref{ch13:calculus} constructs their inclusions
and derivatives, and Theorem~\ref{ch14:rotations} identifies their $O(4)$
representations.

\paragraph{Phases and relativistic fields in (ii).}
Section~\ref{ch15:characters} gives the integral decomposition into phases.
Sections~\ref{ch15:phase-symmetry} and~\ref{ch15:fields} establish symmetry,
reflection positivity and equality of the field spaces in every phase.
Section~\ref{ch16:boundary-section} constructs the tempered boundary
values. Theorem~\ref{ch16:fields} reconstructs the local fields and proves
the spectrum condition. Theorem~\ref{ch16:vacuum} proves vacuum uniqueness
in every phase, and Proposition~\ref{ch16:irreducible} determines the
weak commutant.

\paragraph{The stress tensor and local products in (iii).}
Proposition~\ref{ch14:normalized-stress} fixes the normalization of the
conserved tensor by its Euclidean Ward identities. Section~\ref{ch16:charges-section}
continues them to the relativistic charge law~\eqref{ch16:charge-law}.
Theorem~\ref{ch18:ope-theorem} gives the finite-order expansion and its
remainder estimate. The distinct assertion about extension across the
collision diagonals is Theorem~\ref{ch18:extension-theorem}.

\paragraph{A further conclusion in one selected phase.}
Let $\Theta$ be the trace of the Euclidean stress tensor minus its
expectation in the phase selected in Section~\ref{ch15:selection}, and
let $G_\Theta(r)$ be its two-point function at separation $r>0$ in that
phase. The additional estimates of Chapter~\ref{ch17:trace} give
\[
 \lim_{r\downarrow0} r^8\log^2(1/r)\,G_\Theta(r)
     =\frac{3\dim_{\RR}\gLie}{\pi^4}.
\]
Theorem~\ref{ch17:all-radii} proves the limit for arbitrary radii tending
to zero. Section~\ref{ch17:exclusion} specifies the finite class of
mass-shell correlation functions excluded by this asymptotic.
